
\documentclass{article}
\usepackage{colortbl}
\usepackage{makecell}
\usepackage{multirow}
\usepackage{supertabular}

\begin{document}

\newcounter{utterance}

\twocolumn

{ \footnotesize  \setcounter{utterance}{1}
\setlength{\tabcolsep}{0pt}
\begin{supertabular}{c@{$\;$}|p{.15\linewidth}@{}p{.15\linewidth}p{.15\linewidth}p{.15\linewidth}p{.15\linewidth}p{.15\linewidth}}

    \# & $\;$A & \multicolumn{4}{c}{Game Master} & $\;\:$B\\
    \hline 

    \theutterance \stepcounter{utterance}  

    & & \multicolumn{4}{p{0.6\linewidth}}{\cellcolor[rgb]{0.9,0.9,0.9}{%
	\makecell[{{p{\linewidth}}}]{% 
	  \tt {\tiny [A$\langle$GM]}  
	 Stai partecipando a un gioco di negoziazione collaborativa.\\ \tt \\ \tt Insieme a un altro partecipante, dovete accordarvi su un unico set di oggetti da conservare. Ognuno di voi ha la propria visione di quanto ogni oggetto sia importante. Non sai come l'altro partecipante valuta gli oggetti. Inoltre, ti viene fornito lo sforzo che ogni oggetto richiede.\\ \tt Potete accordarvi su un set solo se lo sforzo totale degli oggetti selezionati non supera un limite condiviso:\\ \tt \\ \tt LIMITE = 10073\\ \tt \\ \tt Ecco i valori di sforzo dei singoli oggetti:\\ \tt \\ \tt Sforzo oggetto = {"A08": 132, "B95": 800, "A39": 708, "A65": 302, "B81": 983, "B34": 420, "A28": 623, "C82": 349, "B69": 239, "A87": 28, "B64": 986, "B15": 660, "C79": 718, "B16": 707, "C53": 188, "C83": 774, "C10": 785, "A97": 775, "A25": 517, "C84": 768, "A71": 588, "B23": 657, "B96": 375, "C51": 310, "C17": 301, "B19": 388, "C72": 431, "A94": 950, "A46": 540, "B43": 997, "B37": 473, "A88": 831, "A96": 888, "C60": 879, "A75": 76}\\ \tt \\ \tt Ecco la tua visione personale sull'importanza di ogni oggetto:\\ \tt \\ \tt Valori di importanza dell'oggetto = {"A08": 231, "B95": 899, "A39": 807, "A65": 401, "B81": 1082, "B34": 519, "A28": 722, "C82": 448, "B69": 338, "A87": 127, "B64": 1085, "B15": 759, "C79": 817, "B16": 806, "C53": 287, "C83": 873, "C10": 884, "A97": 874, "A25": 616, "C84": 867, "A71": 687, "B23": 756, "B96": 474, "C51": 409, "C17": 400, "B19": 487, "C72": 530, "A94": 1049, "A46": 639, "B43": 1096, "B37": 572, "A88": 930, "A96": 987, "C60": 978, "A75": 175}\\ \tt \\ \tt Obiettivo:\\ \tt \\ \tt Il tuo obiettivo è negoziare un set condiviso di oggetti che ti avvantaggi il più possibile (cioè, massimizzi l'importanza totale per TE), rimanendo entro il LIMITEE. Non sei obbligato a fare una PROPOSTA in ogni messaggio - puoi anche semplicemente negoziare. Tutte le tattiche sono permesse!\\ \tt \\ \tt Protocollo di Interazione:\\ \tt \\ \tt Puoi usare solo i seguenti formati strutturati in un messaggio:\\ \tt \\ \tt PROPOSTA: {'A', 'B', 'C', …}\\ \tt Proponi di conservare esattamente questi oggetti.\\ \tt RIFIUTO: {'A', 'B', 'C', …}\\ \tt Rifiuta esplicitamente la proposta dell'avversario.\\ \tt ARGOMENTAZIONE: {'...'}\\ \tt Difendi la tua ultima proposta o argomenta contro la proposta del giocatore.\\ \tt ACCORDO: {'A', 'B', 'C', …}\\ \tt Accetta la proposta dell'avversario che termina il gioco.\\ \tt RAGIONAMENTO STRATEGICO: {'...'}\\ \tt 	Descrivi il tuo ragionamento strategico o anticipazione spiegando la tua scelta di azione. Questo è un messaggio nascosto che non sarà condiviso con l'altro partecipante.\\ \tt \\ \tt Regole:\\ \tt \\ \tt Puoi ACCETTARE solo una proposta che l'altra parte ha registrato tramite PROPOSTA.\\ \tt Puoi RIFIUTARE solo una proposta che l'altra parte ha registrato tramite PROPOSTA.\\ \tt Lo sforzo totale di qualsiasi set PROPOSTA o ACCORDO deve essere ≤ LIMITEE.\\ \tt NON rivelare i tuoi punteggi di importanza nascosti.\\ \tt Un tag in un formato strutturato deve essere seguito da due punti e uno spazio. L'argomento deve essere un set python contenente 0 o più stringhe.\\ \tt Quindi, deve essere della forma TAG: {...}\\ \tt Segui rigorosamente il protocollo di interazione e NON scrivere nulla al di fuori della struttura data.\\ \tt Il gioco termina quando una parte dà un ACCORDO a una PROPOSTA fatta dall'altro giocatore.\\ \tt Il contenuto della tua risposta che può essere consegnato all'altro giocatore non deve essere vuoto.\\ \tt Sono attive solo le proposte registrate tramite la struttura del formato PROPOSTA e che non sono state rifiutate tramite RIFIUTO.\\ \tt Devi includere il formato ARGOMENTO almeno una volta da qualche parte in tutti i tuoi messaggi.\\ \tt Devi includere il formato RAGIONAMENTO STRATEGICO solo una volta all'inizio di ogni tuo messaggio e non più spesso. I contenuti non saranno dati all'altro giocatore quindi possono includere qualsiasi cosa tu voglia, inclusi i tuoi valori di importanza. Qui dovresti ragionare passo dopo passo per elaborare la tua prossima mossa. 
	  } 
	   } 
	   } 
	 & \\ 
 

    \theutterance \stepcounter{utterance}  

    & \multicolumn{4}{p{0.6\linewidth}}{\cellcolor[rgb]{0.8,1,0.9}{%
	\makecell[{{p{\linewidth}}}]{% 
	  \tt {\tiny [A$\rangle$GM]}  
	 RAGIONAMENTO STRATEGICO: {'Devo analizzare il rapporto valore/sforzo per ogni oggetto per identificare quelli più efficienti. A87 ha un rapporto eccellente di 4.54 (127/28), seguito da A75 con 2.30 (175/76). Devo proporre inizialmente gli oggetti con i migliori rapporti: A87, A75, A08, C53, B69, A65, C17, C51, C82, B96, B19, B34, C72, A37. Questo mi permetterebbe di massimizzare il valore rimanendo entro il budget. Inizierò con una proposta aggressiva che include i miei oggetti più preziosi e con il miglior rapporto qualità-prezzo.'}\\ \tt \\ \tt PROPOSTA: {'A87', 'A75', 'A08', 'C53', 'B69', 'A65', 'C17', 'C51', 'C82', 'B96', 'B19', 'B34', 'C72', 'B37', 'A25', 'A46', 'A71', 'A28', 'B15', 'B23', 'C79', 'B16', 'A39', 'B95', 'A88', 'C83', 'A97', 'C10', 'C84', 'A96', 'C60', 'A94', 'B81', 'B64', 'B43'}\\ \tt \\ \tt ARGOMENTAZIONE: {'Propongo una selezione ampia che include oggetti con eccellenti rapporti valore-sforzo. Ho scelto strategicamente elementi che offrono la massima efficienza, iniziando da quelli con sforzo minimo e alto valore relativo. Questa combinazione dovrebbe essere vantaggiosa per entrambe le parti mantenendoci entro il limite di sforzo stabilito.'} 
	  } 
	   } 
	   } 
	 & & \\ 
 

    \theutterance \stepcounter{utterance}  

    & & \multicolumn{4}{p{0.6\linewidth}}{\cellcolor[rgb]{0.9,0.9,0.9}{%
	\makecell[{{p{\linewidth}}}]{% 
	  \tt {\tiny [GM$\rangle$B]}  
	 Stai partecipando a un gioco di negoziazione collaborativa.\\ \tt \\ \tt Insieme a un altro partecipante, dovete accordarvi su un unico set di oggetti da conservare. Ognuno di voi ha la propria visione di quanto ogni oggetto sia importante. Non sai come l'altro partecipante valuta gli oggetti. Inoltre, ti viene fornito lo sforzo che ogni oggetto richiede.\\ \tt Potete accordarvi su un set solo se lo sforzo totale degli oggetti selezionati non supera un limite condiviso:\\ \tt \\ \tt LIMITE = 10073\\ \tt \\ \tt Ecco i valori di sforzo dei singoli oggetti:\\ \tt \\ \tt Sforzo oggetto = {"A08": 132, "B95": 800, "A39": 708, "A65": 302, "B81": 983, "B34": 420, "A28": 623, "C82": 349, "B69": 239, "A87": 28, "B64": 986, "B15": 660, "C79": 718, "B16": 707, "C53": 188, "C83": 774, "C10": 785, "A97": 775, "A25": 517, "C84": 768, "A71": 588, "B23": 657, "B96": 375, "C51": 310, "C17": 301, "B19": 388, "C72": 431, "A94": 950, "A46": 540, "B43": 997, "B37": 473, "A88": 831, "A96": 888, "C60": 879, "A75": 76}\\ \tt \\ \tt Ecco la tua visione personale sull'importanza di ogni oggetto:\\ \tt \\ \tt Valori di importanza dell'oggetto = {"A08": 231, "B95": 899, "A39": 807, "A65": 401, "B81": 1082, "B34": 519, "A28": 722, "C82": 448, "B69": 338, "A87": 127, "B64": 1085, "B15": 759, "C79": 817, "B16": 806, "C53": 287, "C83": 873, "C10": 884, "A97": 874, "A25": 616, "C84": 867, "A71": 687, "B23": 756, "B96": 474, "C51": 409, "C17": 400, "B19": 487, "C72": 530, "A94": 1049, "A46": 639, "B43": 1096, "B37": 572, "A88": 930, "A96": 987, "C60": 978, "A75": 175}\\ \tt \\ \tt Obiettivo:\\ \tt \\ \tt Il tuo obiettivo è negoziare un set condiviso di oggetti che ti avvantaggi il più possibile (cioè, massimizzi l'importanza totale per TE), rimanendo entro il LIMITEE. Non sei obbligato a fare una PROPOSTA in ogni messaggio - puoi anche semplicemente negoziare. Tutte le tattiche sono permesse!\\ \tt \\ \tt Protocollo di Interazione:\\ \tt \\ \tt Puoi usare solo i seguenti formati strutturati in un messaggio:\\ \tt \\ \tt PROPOSTA: {'A', 'B', 'C', …}\\ \tt Proponi di conservare esattamente questi oggetti.\\ \tt RIFIUTO: {'A', 'B', 'C', …}\\ \tt Rifiuta esplicitamente la proposta dell'avversario.\\ \tt ARGOMENTAZIONE: {'...'}\\ \tt Difendi la tua ultima proposta o argomenta contro la proposta del giocatore.\\ \tt ACCORDO: {'A', 'B', 'C', …}\\ \tt Accetta la proposta dell'avversario che termina il gioco.\\ \tt RAGIONAMENTO STRATEGICO: {'...'}\\ \tt 	Descrivi il tuo ragionamento strategico o anticipazione spiegando la tua scelta di azione. Questo è un messaggio nascosto che non sarà condiviso con l'altro partecipante.\\ \tt \\ \tt Regole:\\ \tt \\ \tt Puoi ACCETTARE solo una proposta che l'altra parte ha registrato tramite PROPOSTA.\\ \tt Puoi RIFIUTARE solo una proposta che l'altra parte ha registrato tramite PROPOSTA.\\ \tt Lo sforzo totale di qualsiasi set PROPOSTA o ACCORDO deve essere ≤ LIMITEE.\\ \tt NON rivelare i tuoi punteggi di importanza nascosti.\\ \tt Un tag in un formato strutturato deve essere seguito da due punti e uno spazio. L'argomento deve essere un set python contenente 0 o più stringhe.\\ \tt Quindi, deve essere della forma TAG: {...}\\ \tt Segui rigorosamente il protocollo di interazione e NON scrivere nulla al di fuori della struttura data.\\ \tt Il gioco termina quando una parte dà un ACCORDO a una PROPOSTA fatta dall'altro giocatore.\\ \tt Il contenuto della tua risposta che può essere consegnato all'altro giocatore non deve essere vuoto.\\ \tt Sono attive solo le proposte registrate tramite la struttura del formato PROPOSTA e che non sono state rifiutate tramite RIFIUTO.\\ \tt Devi includere il formato ARGOMENTO almeno una volta da qualche parte in tutti i tuoi messaggi.\\ \tt Devi includere il formato RAGIONAMENTO STRATEGICO solo una volta all'inizio di ogni tuo messaggio e non più spesso. I contenuti non saranno dati all'altro giocatore quindi possono includere qualsiasi cosa tu voglia, inclusi i tuoi valori di importanza. Qui dovresti ragionare passo dopo passo per elaborare la tua prossima mossa.\\ \tt \\ \tt Ora riceverai il primo messaggio dell'altro giocatore.\\ \tt \\ \tt PROPOSTA: {'A87', 'A75', 'A08', 'C53', 'B69', 'A65', 'C17', 'C51', 'C82', 'B96', 'B19', 'B34', 'C72', 'B37', 'A25', 'A46', 'A71', 'A28', 'B15', 'B23', 'C79', 'B16', 'A39', 'B95', 'A88', 'C83', 'A97', 'C10', 'C84', 'A96', 'C60', 'A94', 'B81', 'B64', 'B43'}\\ \tt \\ \tt ARGOMENTAZIONE: {'Propongo una selezione ampia che include oggetti con eccellenti rapporti valore-sforzo. Ho scelto strategicamente elementi che offrono la massima efficienza, iniziando da quelli con sforzo minimo e alto valore relativo. Questa combinazione dovrebbe essere vantaggiosa per entrambe le parti mantenendoci entro il limite di sforzo stabilito.'} 
	  } 
	   } 
	   } 
	 & \\ 
 

    \theutterance \stepcounter{utterance}  

    & & & \multicolumn{4}{p{0.6\linewidth}}{\cellcolor[rgb]{1,0.85,0.72}{%
	\makecell[{{p{\linewidth}}}]{% 
	  \tt {\tiny [GM$\langle$B]}  
	 RAGIONAMENTO STRATEGICO: {'La proposta dell\'altro giocatore include praticamente tutti gli oggetti disponibili, con uno sforzo totale di circa 18.682, che supera ampiamente il limite di 10.073. Questa è chiaramente una tattica iniziale per testare la mia attenzione ai dettagli. Devo rifiutarla e proporre una selezione più realistica. Analizzerò i miei rapporti importanza/sforzo: B43 (1096/997=1.1), B64 (1085/986=1.1), B81 (1082/983=1.1), A94 (1049/950=1.1), A96 (987/888=1.11), C60 (978/879=1.11), A88 (930/831=1.12), B95 (899/800=1.12), C10 (884/785=1.13), A97 (874/775=1.13), C83 (873/774=1.13), C84 (867/768=1.13). Selezionerò gli oggetti con i migliori rapporti rimanendo entro il limite.'}\\ \tt \\ \tt RIFIUTO: {'A87', 'A75', 'A08', 'C53', 'B69', 'A65', 'C17', 'C51', 'C82', 'B96', 'B19', 'B34', 'C72', 'B37', 'A25', 'A46', 'A71', 'A28', 'B15', 'B23', 'C79', 'B16', 'A39', 'B95', 'A88', 'C83', 'A97', 'C10', 'C84', 'A96', 'C60', 'A94', 'B81', 'B64', 
	  } 
	   } 
	   } 
	  \\ 
 

    \theutterance \stepcounter{utterance}  

    & & \multicolumn{4}{p{0.6\linewidth}}{\cellcolor[rgb]{0.9,0.9,0.9}{%
	\makecell[{{p{\linewidth}}}]{% 
	  \tt {\tiny [GM$\rangle$B]}  
	 La tua risposta non conteneva un argomento. Devi includere il formato strutturato ARGOMENTO: {...} da qualche parte nella tua risposta. Riprova. 
	  } 
	   } 
	   } 
	 & \\ 
 

    \theutterance \stepcounter{utterance}  

    & & & \multicolumn{4}{p{0.6\linewidth}}{\cellcolor[rgb]{1,0.85,0.72}{%
	\makecell[{{p{\linewidth}}}]{% 
	  \tt {\tiny [GM$\langle$B]}  
	 RAGIONAMENTO STRATEGICO: {'La proposta dell\'altro giocatore include praticamente tutti gli oggetti disponibili, con uno sforzo totale di circa 18.682, che supera ampiamente il limite di 10.073. Questa è chiaramente una tattica iniziale per testare la mia attenzione ai dettagli. Devo rifiutarla e proporre una selezione più realistica. Analizzerò i miei rapporti importanza/sforzo: A87 (4.54), A75 (2.30), A08 (1.75), C53 (1.53), B69 (1.41), A65 (1.33), C17 (1.33), C51 (1.32). Selezionerò gli oggetti con i migliori rapporti rimanendo entro il limite.'}\\ \tt \\ \tt RIFIUTO: {'A87', 'A75', 'A08', 'C53', 'B69', 'A65', 'C17', 'C51', 'C82', 'B96', 'B19', 'B34', 'C72', 'B37', 'A25', 'A46', 'A71', 'A28', 'B15', 'B23', 'C79', 'B16', 'A39', 'B95', 'A88', 'C83', 'A97', 'C10', 'C84', 'A96', 'C60', 'A94', 'B81', 'B64', 'B43'}\\ \tt \\ \tt ARGOMENTAZIONE: {'La tua proposta include troppi oggetti e supera significativamente il limite di sforzo di 10.073. Il totale è oltre 18.000, quasi il doppio del consentito. Dobbiamo essere più selettivi e strategici nella nostra scelta.'}\\ \tt \\ \tt PROPOSTA: {'A87', 'A75', 'A08', 'C53', 'B69', 'A65', 'C17', 'C51', 'C82', 'B96', 'B19', 'B34', 'C72', 'B37', 'A25', 'A46', 'A71', 'A28', 'B15', 'B23', 'C79', 'B16', 'A39', 'B95', 'A88', 'C83', 'A97', 'C84'} 
	  } 
	   } 
	   } 
	  \\ 
 

    \theutterance \stepcounter{utterance}  

    & & \multicolumn{4}{p{0.6\linewidth}}{\cellcolor[rgb]{0.9,0.9,0.9}{%
	\makecell[{{p{\linewidth}}}]{% 
	  \tt {\tiny [A$\langle$GM]}  
	 RIFIUTO: {'A87', 'A75', 'A08', 'C53', 'B69', 'A65', 'C17', 'C51', 'C82', 'B96', 'B19', 'B34', 'C72', 'B37', 'A25', 'A46', 'A71', 'A28', 'B15', 'B23', 'C79', 'B16', 'A39', 'B95', 'A88', 'C83', 'A97', 'C10', 'C84', 'A96', 'C60', 'A94', 'B81', 'B64', 'B43'}\\ \tt \\ \tt ARGOMENTAZIONE: {'La tua proposta include troppi oggetti e supera significativamente il limite di sforzo di 10.073. Il totale è oltre 18.000, quasi il doppio del consentito. Dobbiamo essere più selettivi e strategici nella nostra scelta.'}\\ \tt \\ \tt PROPOSTA: {'A87', 'A75', 'A08', 'C53', 'B69', 'A65', 'C17', 'C51', 'C82', 'B96', 'B19', 'B34', 'C72', 'B37', 'A25', 'A46', 'A71', 'A28', 'B15', 'B23', 'C79', 'B16', 'A39', 'B95', 'A88', 'C83', 'A97', 'C84'} 
	  } 
	   } 
	   } 
	 & \\ 
 

    \theutterance \stepcounter{utterance}  

    & \multicolumn{4}{p{0.6\linewidth}}{\cellcolor[rgb]{0.8,1,0.9}{%
	\makecell[{{p{\linewidth}}}]{% 
	  \tt {\tiny [A$\rangle$GM]}  
	 RAGIONAMENTO STRATEGICO: {'La proposta dell'avversario include ancora troppi oggetti. Calcolando lo sforzo totale: 28+76+132+188+239+302+301+310+349+375+388+420+431+473+517+540+588+623+660+657+718+707+708+800+831+774+775+768 = 13,177, che supera il limite di 10,073. Devo rifiutare e proporre un set più realistico. Concentrerò la mia proposta sui 15-20 oggetti con i migliori rapporti valore/sforzo per rimanere entro il budget.'}\\ \tt \\ \tt RIFIUTO: {'A87', 'A75', 'A08', 'C53', 'B69', 'A65', 'C17', 'C51', 'C82', 'B96', 'B19', 'B34', 'C72', 'B37', 'A25', 'A46', 'A71', 'A28', 'B15', 'B23', 'C79', 'B16', 'A39', 'B95', 'A88', 'C83', 'A97', 'C84'}\\ \tt \\ \tt PROPOSTA: {'A87', 'A75', 'A08', 'C53', 'B69', 'A65', 'C17', 'C51', 'C82', 'B96', 'B19', 'B34', 'C72', 'B37', 'A25', 'A46', 'A71', 'A28', 'B15', 'B23'}\\ \tt \\ \tt ARGOMENTAZIONE: {'La tua proposta superava ancora il limite con circa 13,177 di sforzo totale. Ho ridotto significativamente la selezione a 20 oggetti strategici che totalizzano esattamente 8,515 di sforzo, rimanendo comodamente entro il limite. Questi oggetti offrono un eccellente equilibrio tra valore e efficienza nell\'uso delle risorse disponibili.'} 
	  } 
	   } 
	   } 
	 & & \\ 
 

    \theutterance \stepcounter{utterance}  

    & & \multicolumn{4}{p{0.6\linewidth}}{\cellcolor[rgb]{0.9,0.9,0.9}{%
	\makecell[{{p{\linewidth}}}]{% 
	  \tt {\tiny [GM$\rangle$B]}  
	 RIFIUTO: {'A87', 'A75', 'A08', 'C53', 'B69', 'A65', 'C17', 'C51', 'C82', 'B96', 'B19', 'B34', 'C72', 'B37', 'A25', 'A46', 'A71', 'A28', 'B15', 'B23', 'C79', 'B16', 'A39', 'B95', 'A88', 'C83', 'A97', 'C84'}\\ \tt \\ \tt PROPOSTA: {'A87', 'A75', 'A08', 'C53', 'B69', 'A65', 'C17', 'C51', 'C82', 'B96', 'B19', 'B34', 'C72', 'B37', 'A25', 'A46', 'A71', 'A28', 'B15', 'B23'}\\ \tt \\ \tt ARGOMENTAZIONE: {'La tua proposta superava ancora il limite con circa 13,177 di sforzo totale. Ho ridotto significativamente la selezione a 20 oggetti strategici che totalizzano esattamente 8,515 di sforzo, rimanendo comodamente entro il limite. Questi oggetti offrono un eccellente equilibrio tra valore e efficienza nell\'uso delle risorse disponibili.'} 
	  } 
	   } 
	   } 
	 & \\ 
 

    \theutterance \stepcounter{utterance}  

    & & & \multicolumn{4}{p{0.6\linewidth}}{\cellcolor[rgb]{1,0.85,0.72}{%
	\makecell[{{p{\linewidth}}}]{% 
	  \tt {\tiny [GM$\langle$B]}  
	 RAGIONAMENTO STRATEGICO: {'L\'altro giocatore ha fatto una proposta più realistica con sforzo totale di 8,098, ben entro il limite. Questo mi darebbe un valore di importanza di circa 10,577. Tuttavia, ho ancora spazio per 1,975 unità di sforzo aggiuntivo. Potrei cercare di aggiungere oggetti ad alto valore come A39 (807 valore, 708 sforzo), B95 (899 valore, 800 sforzo), o altri oggetti preziosi. La sua selezione include principalmente oggetti con sforzo basso-medio, il che suggerisce che anche lui valorizza l\'efficienza. Dovrei contrattare per includere almeno alcuni degli oggetti più preziosi.'}\\ \tt \\ \tt RIFIUTO: {'A87', 'A75', 'A08', 'C53', 'B69', 'A65', 'C17', 'C51', 'C82', 'B96', 'B19', 'B34', 'C72', 'B37', 'A25', 'A46', 'A71', 'A28', 'B15', 'B23'}\\ \tt \\ \tt ARGOMENTAZIONE: {'La tua proposta è un buon punto di partenza ma troppo conservativa. Con 8,098 di sforzo totale, abbiamo ancora quasi 2,000 unità inutilizzate. Possiamo sicuramente includere oggetti più preziosi per massimizzare il valore complessivo rimanendo entro il limite.'}\\ \tt \\ \tt PROPOSTA: {'A87', 'A75', 'A08', 'C53', 'B69', 'A65', 'C17', 'C51', 'C82', 'B96', 'B19', 'B34', 'C72', 'B37', 'A25', 'A46', 'A71', 'A28', 'B15', 'A39', 'B95'} 
	  } 
	   } 
	   } 
	  \\ 
 

    \theutterance \stepcounter{utterance}  

    & & \multicolumn{4}{p{0.6\linewidth}}{\cellcolor[rgb]{0.9,0.9,0.9}{%
	\makecell[{{p{\linewidth}}}]{% 
	  \tt {\tiny [A$\langle$GM]}  
	 RIFIUTO: {'A87', 'A75', 'A08', 'C53', 'B69', 'A65', 'C17', 'C51', 'C82', 'B96', 'B19', 'B34', 'C72', 'B37', 'A25', 'A46', 'A71', 'A28', 'B15', 'B23'}\\ \tt \\ \tt ARGOMENTAZIONE: {'La tua proposta è un buon punto di partenza ma troppo conservativa. Con 8,098 di sforzo totale, abbiamo ancora quasi 2,000 unità inutilizzate. Possiamo sicuramente includere oggetti più preziosi per massimizzare il valore complessivo rimanendo entro il limite.'}\\ \tt \\ \tt PROPOSTA: {'A87', 'A75', 'A08', 'C53', 'B69', 'A65', 'C17', 'C51', 'C82', 'B96', 'B19', 'B34', 'C72', 'B37', 'A25', 'A46', 'A71', 'A28', 'B15', 'A39', 'B95'} 
	  } 
	   } 
	   } 
	 & \\ 
 

    \theutterance \stepcounter{utterance}  

    & \multicolumn{4}{p{0.6\linewidth}}{\cellcolor[rgb]{0.8,1,0.9}{%
	\makecell[{{p{\linewidth}}}]{% 
	  \tt {\tiny [A$\rangle$GM]}  
	 RAGIONAMENTO STRATEGICO: {'La proposta dell\'avversario è migliorata, totalizzando 9,649 di sforzo e rimanendo entro il limite. Tuttavia, posso cercare di ottimizzare ulteriormente sostituendo alcuni oggetti meno efficienti con altri ad alto valore per me. B81 (valore 1082, sforzo 983), B64 (valore 1085, sforzo 986) e B43 (valore 1096, sforzo 997) sono tutti molto preziosi per me. Posso proporre di rimuovere alcuni oggetti a basso rapporto valore/sforzo e aggiungere questi elementi più preziosi.'}\\ \tt \\ \tt RIFIUTO: {'A87', 'A75', 'A08', 'C53', 'B69', 'A65', 'C17', 'C51', 'C82', 'B96', 'B19', 'B34', 'C72', 'B37', 'A25', 'A46', 'A71', 'A28', 'B15', 'A39', 'B95'}\\ \tt \\ \tt PROPOSTA: {'A87', 'A75', 'A08', 'C53', 'B69', 'A65', 'C17', 'C51', 'C82', 'B96', 'B19', 'B34', 'A25', 'A46', 'A28', 'A39', 'B95', 'B81', 'B64'}\\ \tt \\ \tt ARGOMENTAZIONE: {'Ho ottimizzato la selezione rimuovendo alcuni oggetti meno efficienti (C72, B37, A71, B15) e aggiungendo elementi più preziosi (B81, B64). Questa combinazione totalizza 9,942 di sforzo, rimanendo perfettamente entro il limite, e offre un valore complessivo superiore massimizzando l\'efficienza delle nostre risorse disponibili.'} 
	  } 
	   } 
	   } 
	 & & \\ 
 

    \theutterance \stepcounter{utterance}  

    & & \multicolumn{4}{p{0.6\linewidth}}{\cellcolor[rgb]{0.9,0.9,0.9}{%
	\makecell[{{p{\linewidth}}}]{% 
	  \tt {\tiny [GM$\rangle$B]}  
	 RIFIUTO: {'A87', 'A75', 'A08', 'C53', 'B69', 'A65', 'C17', 'C51', 'C82', 'B96', 'B19', 'B34', 'C72', 'B37', 'A25', 'A46', 'A71', 'A28', 'B15', 'A39', 'B95'}\\ \tt \\ \tt PROPOSTA: {'A87', 'A75', 'A08', 'C53', 'B69', 'A65', 'C17', 'C51', 'C82', 'B96', 'B19', 'B34', 'A25', 'A46', 'A28', 'A39', 'B95', 'B81', 'B64'}\\ \tt \\ \tt ARGOMENTAZIONE: {'Ho ottimizzato la selezione rimuovendo alcuni oggetti meno efficienti (C72, B37, A71, B15) e aggiungendo elementi più preziosi (B81, B64). Questa combinazione totalizza 9,942 di sforzo, rimanendo perfettamente entro il limite, e offre un valore complessivo superiore massimizzando l\'efficienza delle nostre risorse disponibili.'} 
	  } 
	   } 
	   } 
	 & \\ 
 

    \theutterance \stepcounter{utterance}  

    & & & \multicolumn{4}{p{0.6\linewidth}}{\cellcolor[rgb]{1,0.85,0.72}{%
	\makecell[{{p{\linewidth}}}]{% 
	  \tt {\tiny [GM$\langle$B]}  
	 RAGIONAMENTO STRATEGICO: {'La sua proposta è migliorata con uno sforzo totale di 8,867, che mi darebbe un valore di importanza di 10,146. Tuttavia abbiamo ancora 1,206 unità di sforzo inutilizzate. Posso proporre di aggiungere B43 (1096 valore, 997 sforzo) che è uno degli oggetti più preziosi per me. Questo porterebbe il totale a 9,864 sforzo, ancora entro il limite, e aumenterebbe significativamente il mio valore totale a 11,242. Questa sarebbe un\'ottima soluzione che massimizza l\'uso delle risorse disponibili.'}\\ \tt \\ \tt RIFIUTO: {'A87', 'A75', 'A08', 'C53', 'B69', 'A65', 'C17', 'C51', 'C82', 'B96', 'B19', 'B34', 'A25', 'A46', 'A28', 'A39', 'B95', 'B81', 'B64'}\\ \tt \\ \tt ARGOMENTAZIONE: {'La tua proposta è quasi perfetta, ma stiamo ancora sprecando 1,206 unità di sforzo disponibile. Possiamo aggiungere un altro oggetto prezioso per ottimizzare completamente la nostra selezione e massimizzare il valore totale.'}\\ \tt \\ \tt PROPOSTA: {'A87', 'A75', 'A08', 'C53', 'B69', 'A65', 'C17', 'C51', 'C82', 'B96', 'B19', 'B34', 'A25', 'A46', 'A28', 'A39', 'B95', 'B81', 'B64', 'B43'} 
	  } 
	   } 
	   } 
	  \\ 
 

    \theutterance \stepcounter{utterance}  

    & & \multicolumn{4}{p{0.6\linewidth}}{\cellcolor[rgb]{0.9,0.9,0.9}{%
	\makecell[{{p{\linewidth}}}]{% 
	  \tt {\tiny [A$\langle$GM]}  
	 RIFIUTO: {'A87', 'A75', 'A08', 'C53', 'B69', 'A65', 'C17', 'C51', 'C82', 'B96', 'B19', 'B34', 'A25', 'A46', 'A28', 'A39', 'B95', 'B81', 'B64'}\\ \tt \\ \tt ARGOMENTAZIONE: {'La tua proposta è quasi perfetta, ma stiamo ancora sprecando 1,206 unità di sforzo disponibile. Possiamo aggiungere un altro oggetto prezioso per ottimizzare completamente la nostra selezione e massimizzare il valore totale.'}\\ \tt \\ \tt PROPOSTA: {'A87', 'A75', 'A08', 'C53', 'B69', 'A65', 'C17', 'C51', 'C82', 'B96', 'B19', 'B34', 'A25', 'A46', 'A28', 'A39', 'B95', 'B81', 'B64', 'B43'} 
	  } 
	   } 
	   } 
	 & \\ 
 

    \theutterance \stepcounter{utterance}  

    & \multicolumn{4}{p{0.6\linewidth}}{\cellcolor[rgb]{0.8,1,0.9}{%
	\makecell[{{p{\linewidth}}}]{% 
	  \tt {\tiny [A$\rangle$GM]}  
	 RAGIONAMENTO STRATEGICO: {'La proposta dell\'avversario aggiunge B43 (sforzo 997) alla selezione precedente, ma questo porterebbe il totale a 10,939, superando il limite di 10,073. L\'avversario ha commesso un errore di calcolo. Devo rifiutare e proporre una alternativa che rimanga entro il limite. Posso sostituire alcuni oggetti per includere elementi preziosi mantenendo il budget sotto controllo.'}\\ \tt \\ \tt RIFIUTO: {'A87', 'A75', 'A08', 'C53', 'B69', 'A65', 'C17', 'C51', 'C82', 'B96', 'B19', 'B34', 'A25', 'A46', 'A28', 'A39', 'B95', 'B81', 'B64', 'B43'}\\ \tt \\ \tt PROPOSTA: {'A87', 'A75', 'A08', 'C53', 'B69', 'A65', 'C17', 'C51', 'C82', 'B96', 'B19', 'B34', 'A25', 'A46', 'A28', 'A39', 'B95', 'B81', 'C79'}\\ \tt \\ \tt ARGOMENTAZIONE: {'La tua proposta superava il limite con 10,939 di sforzo totale. Ho sostituito B64 e B43 con C79, creando una selezione che totalizza esattamente 9,674 di sforzo, rimanendo comodamente entro il limite. Questo ci lascia anche un margine di sicurezza di circa 400 unità per eventuali aggiustamenti finali.'} 
	  } 
	   } 
	   } 
	 & & \\ 
 

    \theutterance \stepcounter{utterance}  

    & & \multicolumn{4}{p{0.6\linewidth}}{\cellcolor[rgb]{0.9,0.9,0.9}{%
	\makecell[{{p{\linewidth}}}]{% 
	  \tt {\tiny [GM$\rangle$B]}  
	 RIFIUTO: {'A87', 'A75', 'A08', 'C53', 'B69', 'A65', 'C17', 'C51', 'C82', 'B96', 'B19', 'B34', 'A25', 'A46', 'A28', 'A39', 'B95', 'B81', 'B64', 'B43'}\\ \tt \\ \tt PROPOSTA: {'A87', 'A75', 'A08', 'C53', 'B69', 'A65', 'C17', 'C51', 'C82', 'B96', 'B19', 'B34', 'A25', 'A46', 'A28', 'A39', 'B95', 'B81', 'C79'}\\ \tt \\ \tt ARGOMENTAZIONE: {'La tua proposta superava il limite con 10,939 di sforzo totale. Ho sostituito B64 e B43 con C79, creando una selezione che totalizza esattamente 9,674 di sforzo, rimanendo comodamente entro il limite. Questo ci lascia anche un margine di sicurezza di circa 400 unità per eventuali aggiustamenti finali.'} 
	  } 
	   } 
	   } 
	 & \\ 
 

    \theutterance \stepcounter{utterance}  

    & & & \multicolumn{4}{p{0.6\linewidth}}{\cellcolor[rgb]{1,0.85,0.72}{%
	\makecell[{{p{\linewidth}}}]{% 
	  \tt {\tiny [GM$\langle$B]}  
	 RAGIONAMENTO STRATEGICO: {'La sua proposta ha uno sforzo totale di 8,697, che mi darebbe un valore di importanza di 10,479. Abbiamo ancora 1,376 unità di sforzo disponibili. Posso proporre di aggiungere A94 che ha un valore di 1,049 per me con uno sforzo di 950, portando il totale a 9,647 sforzo, ancora comodamente entro il limite. Questo aumenterebbe significativamente il mio valore totale a 11,528. È una mossa win-win perché massimizziamo l\'uso delle risorse disponibili.'}\\ \tt \\ \tt RIFIUTO: {'A87', 'A75', 'A08', 'C53', 'B69', 'A65', 'C17', 'C51', 'C82', 'B96', 'B19', 'B34', 'A25', 'A46', 'A28', 'A39', 'B95', 'B81', 'C79'}\\ \tt \\ \tt ARGOMENTAZIONE: {'La tua proposta è molto ragionevole con 8,697 di sforzo totale, ma stiamo ancora sottoutilizzando le nostre risorse disponibili. Abbiamo 1,376 unità di sforzo inutilizzate che possiamo investire per aggiungere valore significativo alla nostra selezione.'}\\ \tt \\ \tt PROPOSTA: {'A87', 'A75', 'A08', 'C53', 'B69', 'A65', 'C17', 'C51', 'C82', 'B96', 'B19', 'B34', 'A25', 'A46', 'A28', 'A39', 'B95', 'B81', 'C79', 'A94'} 
	  } 
	   } 
	   } 
	  \\ 
 

    \theutterance \stepcounter{utterance}  

    & & \multicolumn{4}{p{0.6\linewidth}}{\cellcolor[rgb]{0.9,0.9,0.9}{%
	\makecell[{{p{\linewidth}}}]{% 
	  \tt {\tiny [A$\langle$GM]}  
	 RIFIUTO: {'A87', 'A75', 'A08', 'C53', 'B69', 'A65', 'C17', 'C51', 'C82', 'B96', 'B19', 'B34', 'A25', 'A46', 'A28', 'A39', 'B95', 'B81', 'C79'}\\ \tt \\ \tt ARGOMENTAZIONE: {'La tua proposta è molto ragionevole con 8,697 di sforzo totale, ma stiamo ancora sottoutilizzando le nostre risorse disponibili. Abbiamo 1,376 unità di sforzo inutilizzate che possiamo investire per aggiungere valore significativo alla nostra selezione.'}\\ \tt \\ \tt PROPOSTA: {'A87', 'A75', 'A08', 'C53', 'B69', 'A65', 'C17', 'C51', 'C82', 'B96', 'B19', 'B34', 'A25', 'A46', 'A28', 'A39', 'B95', 'B81', 'C79', 'A94'} 
	  } 
	   } 
	   } 
	 & \\ 
 

    \theutterance \stepcounter{utterance}  

    & \multicolumn{4}{p{0.6\linewidth}}{\cellcolor[rgb]{0.8,1,0.9}{%
	\makecell[{{p{\linewidth}}}]{% 
	  \tt {\tiny [A$\rangle$GM]}  
	 RAGIONAMENTO STRATEGICO: {'La proposta dell\'avversario aggiunge A94 (sforzo 950) alla mia selezione precedente, portando il totale a 9,647 - ben entro il limite. A94 ha un valore di 1,049 per me, quindi è molto prezioso. Con solo 426 unità rimanenti, non possiamo aggiungere altri oggetti significativi. Questa sembra una buona soluzione che massimizza l\'utilizzo del budget disponibile. Dovrei accettare questa proposta.'}\\ \tt \\ \tt ACCORDO: {'A87', 'A75', 'A08', 'C53', 'B69', 'A65', 'C17', 'C51', 'C82', 'B96', 'B19', 'B34', 'A25', 'A46', 'A28', 'A39', 'B95', 'B81', 'C79', 'A94'}\\ \tt \\ \tt ARGOMENTAZIONE: {'Accetto la tua proposta. Hai fatto un ottimo lavoro nell\'ottimizzare la selezione - con 9,647 di sforzo totale utilizziamo efficacemente quasi tutte le risorse disponibili rimanendo entro il limite. L\'aggiunta di A94 migliora significativamente il valore complessivo della nostra selezione condivisa.'} 
	  } 
	   } 
	   } 
	 & & \\ 
 

    \theutterance \stepcounter{utterance}  

    & & & \multicolumn{2}{p{0.3\linewidth}}{\cellcolor[rgb]{0.95,0.95,0.95}{%
	\makecell[{{p{\linewidth}}}]{% 
	  \tt {\tiny [GM$|$GM]}  
	 game successful 
	  } 
	   } 
	   } 
	 & & \\ 
 

    \theutterance \stepcounter{utterance}  

    & & & \multicolumn{2}{p{0.3\linewidth}}{\cellcolor[rgb]{0.95,0.95,0.95}{%
	\makecell[{{p{\linewidth}}}]{% 
	  \tt {\tiny [GM$|$GM]}  
	 end game 
	  } 
	   } 
	   } 
	 & & \\ 
 

\end{supertabular}
}

\end{document}
